\chapter{June Huh}\label{ch2022junehuh:june-huh-fields-medal-2022-a-comprehensive-report-on-his-work-and-impact}

\begin{center}\rule{0.5\linewidth}{0.5pt}\end{center}

\textbf{Prepared:} July 2026

\textbf{Subject:} June Huh (b.\\ June 9, 1983, Washington, D.C., USA)

\textbf{Position:} Professor, School of Mathematics, Institute for Advanced Study, Princeton

\textbf{Doctoral Advisor:} Mircea Musta\c{t}\u{a} (University of Michigan)

\textbf{Fields Medal Citation:} \emph{``For bringing the ideas of Hodge theory to combinatorics, the proof of the Dowling--Wilson conjecture for geometric lattices and the Heron--Rota--Welsh conjecture for matroids.''}

\begin{center}\rule{0.5\linewidth}{0.5pt}\end{center}

\begin{prereqs}
Combinatorics (matroids, graphs, independent sets), algebraic geometry (Hodge theory, complex manifolds, Chow rings), basic linear algebra, topology (cohomology of manifolds).
\end{prereqs}

\begin{objectives}
Explain the Heron--Rota--Welsh conjecture and its proof via Hodge theory; describe the intersection theory of matroid Schubert varieties; understand the log-concavity of the characteristic polynomial and its applications.
\end{objectives}

\section{Problem History}

A \textbf{matroid} is a combinatorial abstraction of the notion of independence, unifying linear algebra (independent vectors), graph theory (forests), and field theory (algebraic independence). Given a matroid $M$ on ground set $E$ of rank $r$, the \textbf{characteristic polynomial} is:

\[
\chi_M(q) = \sum_{i=0}^r (-1)^i w_i q^{r-i}
\]

where $w_i$ is the number of independent sets of rank $i$ (the Whitney numbers of the second kind). The \textbf{Heron--Rota--Welsh conjecture} (1977) states that the coefficients of $\chi_M(q)$ are unimodal for any matroid $M$: there exists $j$ such that $w_0 \le w_1 \le \cdots \le w_j \ge w_{j+1} \ge \cdots$.

This is far from obvious. For the uniform matroid $U_{r,n}$ (all subsets of size $\le r$ are independent), the characteristic polynomial is:

\[
\chi_{U_{r,n}}(q) = \sum_{i=0}^{r-1} \binom{n}{i}(-1)^i q^{r-i}
\]

and proving unimodality of binomial coefficients is a classical result. But the conjecture was open for general matroids.

The \textbf{Dowling--Wilson conjecture} (1973) predicted that the Whitney numbers of geometric lattices are unimodal. This was known to follow from the Heron--Rota--Welsh conjecture via the relationship between geometric lattices and simple matroids.

G\'al's conjecture (2010) proposed that the coefficients of the characteristic polynomial of the graphic matroid of a complete graph $K_n$ satisfy a strong form of unimodality (log-concavity) related to the coefficients of the Riemann zeta function.

The theory of \textbf{matroids} was introduced by Whitney (1935) as an abstraction of linear independence. A matroid $M$ on ground set $E$ is a family of subsets (the \textbf{independent sets}) satisfying the axioms of empty set, heredity, and the augmentation property. The rank function $r: 2^E \to \mathbb{Z}_{\ge 0}$ encodes the combinatorial structure, and the characteristic polynomial $\chi_M(q)$ is defined using the Whitney numbers. The question of whether the coefficients of $\chi_M(q)$ exhibit nice properties such as unimodality or log-concavity had been open since the 1970s.

The \textbf{Tutte polynomial} of a matroid $M$ is a two-variable polynomial $T_M(x,y)$ that encodes many combinatorial invariants: the characteristic polynomial is a specialization $T_M(1-q, 0) = (-1)^{r(M)} q^{r(M)} \chi_M(q)$, and the number of bases, independent sets, and other quantities can be extracted from $T_M$. The log-concavity of the characteristic polynomial is equivalent to a log-concavity property of certain specializations of the Tutte polynomial, which connects Huh's work \cite{huh2015} to the broader theory of matroid invariants.

The \textbf{Hodge index theorem} for smooth projective varieties states that the signature of the intersection form on the cohomology ring alternates in sign. The hard Lefschetz property, proved by Deligne, implies that the h-vector of a smooth projective variety is log-concave. Huh's key insight was to find a combinatorial analogue of this property for matroids.

\section{How Was It Solved}

Huh's approach was to interpret matroid combinatorics using algebraic geometry, following an intuition from the theory of mixed Hodge structures. The central idea is to associate to each matroid a geometric object (the Bergman fan) whose cohomological invariants encode the combinatorial data of the matroid.

\textbf{Matroid Schubert varieties.} Huh and Katz introduced the Chow ring of the Bergman fan of a matroid, a combinatorial analogue of the Chow ring of a smooth projective variety. The \textbf{h-vector} of this ring corresponds to the coefficients of the characteristic polynomial. The key insight is that the matroid Chow ring satisfies the hard Lefschetz property, which forces the h-vector to be log-concave.

\textbf{Log-concavity via Hodge theory.} Huh proved the Heron--Rota--Welsh conjecture using Hodge theory. The proof proceeds in three steps: (1) construct a K\"ahler cone on the Bergman fan; (2) prove the hard Lefschetz theorem for the matroid Chow ring using the theory of mixed Hodge structures; (3) deduce log-concavity from the hard Lefschetz property. Log-concavity implies unimodality. The hard Lefschetz property for matroid Chow rings was proved using a degeneration argument: the matroid Chow ring is the associated graded ring of a filtration on the Chow ring of a smooth toric variety, and the hard Lefschetz property is preserved under associated graded rings.

The \textbf{Bergman fan} of a matroid $M$ on ground set $E$ is a subfan of the normal fan of the matroid polytope. It is a pure fan of dimension $r(M)$ (the rank of $M$) whose cones correspond to chains of flats of $M$. The Chow ring of the Bergman fan is defined as the Chow ring of the corresponding toric variety, and its cohomological properties encode the combinatorial data of the matroid.

The \textbf{matroid polytope} $P_M$ is the convex hull of the indicator vectors of the bases of $M$ in $\mathbb{R}^E$. Its face lattice is isomorphic to the lattice of flats of $M$, and the Bergman fan is the normal fan of $P_M$. The Chow ring of the Bergman fan is a quotient of the polynomial ring $\mathbb{Z}[x_e : e \in E]$ by an ideal generated by linear and quadratic relations that encode the combinatorial structure of $M$. The h-vector of this ring is the sequence of Whitney numbers $w_i$, and the hard Lefschetz property implies that these numbers are log-concave.

The connection between the Bergman fan and the matroid polytope provides a geometric interpretation of the log-concavity result: the log-concavity of the Whitney numbers is equivalent to the log-concavity of the volumes of the faces of the matroid polytope. This geometric interpretation suggests connections to the Alexandrov--Fenchel inequality for mixed volumes of convex bodies, which is a deep result in convex geometry with many combinatorial consequences.

\textbf{G\'al's conjecture.} Huh proved G\'al's conjecture for all matroids: the coefficients $w_i$ of the characteristic polynomial satisfy $w_i^2 \ge w_{i-1} w_{i+1}$ for all $i$. The proof uses the log-concavity of the coefficients of the characteristic polynomial, which follows from the hard Lefschetz property of the matroid Chow ring.

G\'al's conjecture was originally formulated for the graphic matroid of the complete graph $K_n$, where the characteristic polynomial is related to the chromatic polynomial. The chromatic polynomial $\chi_{K_n}(q) = q(q-1)(q-2)\cdots(q-n+1)$ counts the number of proper colorings of $K_n$ with $q$ colors, and its coefficients satisfy a strong form of log-concavity. Huh's proof extends this result to all matroids, showing that the log-concavity of the characteristic polynomial is a universal property of matroids that follows from the Hodge-theoretic structure of the matroid Chow ring.

\section{Impact of This Discovery}

The Heron--Rota--Welsh conjecture had been open for over 40 years. Huh's proof established a completely new paradigm: using methods from algebraic geometry and Hodge theory to solve combinatorial problems. This approach has been extended to many other conjectures about matroids.

The Dowling--Wilson conjecture was resolved as a corollary. The connection between matroids and algebraic geometry, established by Huh and Katz, opened new directions in tropical geometry and combinatorics.

The log-concavity results have applications to the theory of stable polynomials (Borcea--Br\"and'en), the theory of real-rooted polynomials, and to understanding the relationship between different combinatorial invariants. The connection to stable polynomials is particularly important: a polynomial with real coefficients is stable if it does not vanish whenever all its variables have positive imaginary part. Stable polynomials are log-concave, and many combinatorial sequences arise as restrictions of stable polynomials.

The \textbf{characteristic polynomial} of a matroid $M$ can be expressed in terms of the Tutte polynomial:
\[
\chi_M(q) = (-1)^{r(M)} q^{c(M)} T_M(1-q, 0),
\]
where $r(M)$ is the rank of $M$ and $c(M)$ is the number of connected components. The Tutte polynomial $T_M(x,y)$ encodes many combinatorial invariants: $T_M(1,1)$ is the number of bases, $T_M(2,1)$ is the number of independent sets, and $T_M(1,2)$ is the number of spanning sets. The log-concavity of the characteristic polynomial is equivalent to a log-concavity property of certain specializations of the Tutte polynomial, which connects Huh's work \cite{huh2015} to the broader theory of matroid invariants.

The \textbf{Alexandrov--Fenchel inequality} for mixed volumes of convex bodies provides the geometric foundation: for convex bodies $K_1, \ldots, K_n$ in $\mathbb{R}^n$, the mixed volumes satisfy a quadratic inequality that implies log-concavity of successive coefficients. Huh's key insight was that the matroid Chow ring satisfies a combinatorial analogue of this inequality, yielding log-concavity of the Whitney numbers.

\section{Modern Views and Improvements}

Huh continues to develop the connection between matroids and algebraic geometry. His recent work on valuations on matroid polyhedra and on the intersection theory of tropical varieties extends the Hodge-theoretic framework.

The theory of matroid Chow rings, developed by Braden--Huh--Matherne--Proudfoot--Wang, provides the foundation for proving new combinatorial identities and inequalities. Extensions to oriented matroids and to other combinatorial structures are active areas.

The methods have connections to representation theory (Kazhdan--Lusztig theory) and to the study of flag manifolds, providing new geometric insights into classical combinatorial objects.

The Kazhdan--Lusztig polynomials of a Coxeter group are a family of polynomials that encode the singularities of Schubert varieties in flag manifolds. The log-concavity of these polynomials is a major open problem, and Huh's methods provide a framework for attacking it. The connection between matroid theory and representation theory has been further developed by Braden, Huh, Matherne, Proudfoot, and Wang, who have established new results on the Kazhdan--Lusztig polynomials of matroids using the Hodge-theoretic framework.

The theory of tropical geometry provides another connection between Huh's work \cite{huh2015} and algebraic geometry. The Bergman fan of a matroid is a tropical variety, and the Chow ring of the Bergman fan is a tropical Chow ring. The hard Lefschetz property for tropical varieties is a major open problem, and Huh's results provide the first examples where it is known to hold. The connection between tropical geometry and matroid theory is an active area of research, with potential applications to enumerative geometry and mirror symmetry.

The \textbf{matroid base polytope} is the convex hull of the indicator vectors of the bases of a matroid. Its geometry is closely related to the combinatorics of the matroid: the faces of the matroid base polytope correspond to the restricted and contracted matroids. The log-concavity of the characteristic polynomial is equivalent to the log-concavity of the volumes of the faces of the matroid base polytope, which is a geometric statement about convex bodies.

Recent work of Adiprasito, Huh, and Katz on the \textbf{Hodge theory for matroids} extends the original results to a broader framework. The key innovation is the construction of a Hodge structure on the cohomology of the Bergman fan, which provides a systematic approach to proving log-concavity and unimodality results. This framework has been applied to prove new inequalities for the Tutte polynomial of a matroid, which encodes many combinatorial invariants.

The connection between Huh's work \cite{huh2015} and the theory of stable polynomials provides another perspective on log-concavity. A polynomial with real coefficients is stable if it does not vanish whenever all its variables have positive imaginary part. The theory of stable polynomials, developed by Borcea and Br\"and'en, provides a systematic framework for proving log-concavity results. The connection between Huh's Hodge-theoretic methods and the theory of stable polynomials is an active area of research.

The theory of matroid valuations, developed by Adiprasito, Huh, and Katz, provides a framework for understanding the log-concavity of the characteristic polynomial through the lens of valuation theory. A valuation on a matroid is a function that satisfies the inclusion-exclusion principle, and the log-concavity of the characteristic polynomial can be viewed as a consequence of the valuation-theoretic properties of the Whitney numbers. This perspective connects Huh's work \cite{huh2015} to the classical theory of valuations in convex geometry and to the theory of Ehrhart polynomials of polytopes.

\begin{center}\rule{0.5\linewidth}{0.5pt}\end{center}

\section*{Summary}
\addcontentsline{toc}{section}{Summary}

Huh proved the Heron--Rota--Welsh conjecture and related results by applying Hodge theory to matroids. His work established deep connections between algebraic geometry and combinatorics, resolving conjectures that had been open for decades.
